\documentclass[border=5mm]{standalone}
% ==============================================================================
% SETUP.TEX - CONFIGURACIÓN GLOBAL DEL PROYECTO
% ==============================================================================
% Este archivo centraliza todos los paquetes y estilos.
% Debe incluirse en el preámbulo del libro principal y de los archivos standalone.
% \PassOptionsToPackage{gray}{xcolor}
% --- 1. CODIFICACIÓN E IDIOMA ---
% fix-cm habilita las fuentes Computer Modern en tamaños arbitrarios (por debajo
% de 5 pt, entre otros). Debe cargarse antes que fontenc.
\usepackage{fix-cm}
\usepackage[utf8]{inputenc}
\usepackage[T1]{fontenc}
\usepackage[spanish, es-tabla]{babel}  
\usepackage{csquotes} % Recomendado para babel y citas

\usepackage{import}
% mode=image: Usa el PDF pre-compilado, nunca intenta recompilar.
% Debes compilar cada figura manualmente antes de compilar el libro.
\usepackage[mode=image, subpreambles=false]{standalone}

% --- 2. MATEMÁTICAS Y CIENCIAS ---
\usepackage{amsmath, amssymb, amsfonts, mathtools}
\usepackage{enumitem}

% LaTeX trae \DeclareMathSizes{5}{5}{5}{5}, así que a 5 pt (\tiny) los subíndices
% salen del mismo tamaño que la base: en las etiquetas de figuras el M de
% $\omega_M$ queda desproporcionado. Se restablece la proporción habitual.
\DeclareMathSizes{5}{5}{3.7}{3}

% --- 3. DISEÑO Y GEOMETRÍA --- 
% Unificamos los márgenes para todo el documento
\usepackage[left=2.5cm,right=2.5cm,top=3cm,bottom=3cm]{geometry}
\makeatletter
\ifstandalone % Evita que geometry dañe el recorte de las figuras bajándolas 3cm
  \@ifclasswith{standalone}{crop=false}{
    % Es un capítulo/sección: Dejamos que geometry actúe.
  }{
    % Es una figura: Neutralizamos geometry para que no las recorte mal.
    \geometry{pass}
  }
\fi
\makeatother
\usepackage{parskip} % Espaciado entre párrafos sin sangría
\usepackage{float}   % Para usar [H] en figura

% Ajustes de flotantes para evitar espacios verticales exagerados
\renewcommand{\topfraction}{0.95}      % Permite que las figuras ocupen hasta el 95% superior
\renewcommand{\bottomfraction}{0.95}   % Permite que las figuras ocupen hasta el 95% inferior
\renewcommand{\textfraction}{0.05}     % Permite hojas con tan solo un 5% de texto
\renewcommand{\floatpagefraction}{0.8} % Requiere que una página de solo figuras esté 80% llena
\setcounter{topnumber}{4}
\setcounter{bottomnumber}{4}
\setcounter{totalnumber}{10}

% --- 4. GRÁFICOS Y TABLAS ---
\usepackage{graphicx}
\usepackage{booktabs} % Tablas profesionales
\usepackage{caption}  % Control de captions y \captionof
\usepackage{tikz}
\usetikzlibrary{arrows.meta, calc, angles, quotes, babel, backgrounds}
\usepackage{pgfplots}
\usepgfplotslibrary{groupplots}
\usepgfplotslibrary{external}
\usepgfplotslibrary{patchplots}
\pgfplotsset{compat=1.18}
\usepackage[american, straightvoltages]{circuitikz}

% --- 5. COLORES Y CAJAS ---

\usepackage[table]{xcolor}
\usepackage{tcolorbox}

% Definición de colores corporativos/estilo
\definecolor{utnblue}{RGB}{0, 51, 102}
\definecolor{codegreen}{rgb}{0,0.6,0}
\definecolor{codegray}{rgb}{0.5,0.5,0.5}
\definecolor{codepurple}{rgb}{0.58,0,0.82}
\definecolor{backcolour}{rgb}{0.96,0.96,0.96}

% --- 6. ENTORNOS PERSONALIZADOS (CAJAS) ---

% Caja "Resultado" (Azul Corporativo) — Definiciones, fórmulas clave, resultados importantes.
% Uso: \begin{resultado}[Fórmula de Euler] ... \end{resultado}
\newtcolorbox{resultado}[1][]{
    colback=utnblue!5!white,
    colframe=utnblue,
    title=\textbf{#1},
    fonttitle=\bfseries,
    sharp corners,
    boxrule=0.5mm
}

% Caja "Nota" (Gris) — Advertencias, aplicaciones, contexto secundario.
% Uso: \begin{nota}[Cuidado con la calculadora] ... \end{nota}
% Uso: \begin{nota} ... \end{nota}  → título por defecto "Nota"
\newtcolorbox{nota}[1][Nota]{
    colback=codegray!5!white,
    colframe=codegray,
    title=\textbf{#1},
    fonttitle=\bfseries,
    sharp corners,
    boxrule=0.5mm
}

% Caja inline para resaltar un valor y enlazarlo con su otra aparición en una tabla
% o en una cuenta: mismo color, mismo valor.
% Uso: \marcavalor{$4$}  o  \marcavalor[codegreen]{$4$}
\newtcbox{\marcavalor}[1][utnblue]{
    on line,
    boxsep=1pt, left=2pt, right=2pt, top=1pt, bottom=1pt,
    colback=#1!10!white,
    colframe=#1,
    boxrule=0.4pt,
    arc=2pt
}

% --- 7. CÓDIGO FUENTE (LISTINGS) ---
\usepackage{listings}

\lstdefinestyle{miestilo}{
    backgroundcolor=\color{backcolour},   
    commentstyle=\color{codegreen},
    keywordstyle=\color{magenta},
    numberstyle=\tiny\color{codegray},
    stringstyle=\color{codepurple},
    basicstyle=\ttfamily\footnotesize,
    breakatwhitespace=false,         
    breaklines=true,                 
    captionpos=b,                    
    keepspaces=true,                 
    numbers=left,                    
    numbersep=5pt,                  
    showspaces=false,                
    showstringspaces=false,
    showtabs=false,                  
    tabsize=2,
    frame=single,                    
    rulecolor=\color{codegray}
}

\lstset{style=miestilo}
% Soporte para tildes y ñ en bloques de código
\lstset{literate={ñ}{{\~n}}1 {á}{{\'a}}1 {é}{{\'e}}1 {í}{{\'i}}1 {ó}{{\'o}}1 {ú}{{\'u}}1}

% Operadores matemáticos personalizados

% Vector en sentido discreto (lista finita de coordenadas): negrita + flecha.
% Uso: \vect{v}, \vect{e}_k, \vect{0}.
\newcommand{\vect}[1]{\vec{\mathbf{#1}}}

% Fecha de versión y comando para capítulos standalone
\providecommand{\verdate}{\the\year.\ifnum\month<10 0\fi\the\month.\ifnum\day<10 0\fi\the\day}
\newcommand{\chapterversion}{%
    \vspace{-1.5cm}\begin{flushright}\small\textit{\color{codegray}Versión~\verdate}\end{flushright}\vspace{0.5cm}%
}

% --- 8. HIPERVÍNCULOS (DEBE SER EL ÚLTIMO PAQUETE) ---
\usepackage[colorlinks=true, linkcolor=utnblue, urlcolor=utnblue, citecolor=utnblue]{hyperref}

% --- 9. REFERENCIAS CRUZADAS ENTRE CAPÍTULOS STANDALONE ---
% Cuando se compila un capítulo standalone, xr-hyper lee los .aux de los
% otros capítulos (si existen) para resolver \ref{} a labels externos.
% En la compilación del libro completo este bloque no se activa.
\ifstandalone
  \usepackage{xr-hyper}
  \makeatletter
  \newcommand{\xrchap}[1]{%
    \edef\xrc@self{\detokenize{Capitulo#1}}%
    \edef\xrc@job{\jobname}%
    \ifx\xrc@self\xrc@job\else
      \IfFileExists{../#1capitulo/Capitulo#1.aux}%
        {\externaldocument{../#1capitulo/Capitulo#1}}{}%
    \fi
  }
  \makeatother
  \xrchap{01}\xrchap{02}\xrchap{03}\xrchap{04}
  \xrchap{05}\xrchap{06}\xrchap{07}\xrchap{08}
  \xrchap{09}\xrchap{10}\xrchap{11}\xrchap{12}
  \xrchap{13}
\fi
% \PassOptionsToPackage{gray}{xcolor}

% fig_periodica_como_impulsos
% Tres paneles lado a lado sobre el tren periodico de pulsos rectangulares.
% Parametros: tau = 1, T_0 = 4, omega_0 = pi/2.
% (a) x(t): tren de pulsos en el tiempo.
% (b) c_k vs k: espectro de lineas (sticks sobre indice discreto).
% (c) X(jw) vs omega: tren de impulsos ponderados por 2*pi*c_k en omega = k*omega_0.
% c_k = (tau/T_0) * sinc(k*omega_0*tau/2) = (1/4) * sinc(k*pi/4).
% 2*pi*c_k = (pi/2)*sinc(k*pi/4).
% Sizing: 3 paneles lado a lado, w 4.5cm, h 4cm.

\begin{document}
\begin{tikzpicture}[>={Latex[length=3pt,width=3pt]}, font=\scriptsize]

% ===== Panel (a): x(t) tren de pulsos =====
\begin{axis}[
    name=panel_a,
    width=4.5cm, height=4cm,
    axis lines=center,
    axis line style={->, thick},
    xmin=-7, xmax=7, ymin=-0.35, ymax=1.55,
    xtick={-4,0,4},
    xticklabels={$-T_0$,$0$,$T_0$},
    ytick={0},
    every tick label/.style={font=\footnotesize},
    xlabel={$t$},
    xlabel style={anchor=west, at={(axis description cs:1.0,0.20)}},
    ylabel={$x(t)$},
    ylabel style={font=\scriptsize, rotate=0, anchor=south east,
                  at={(axis description cs:0,1)}, xshift=-2pt},
    title={(a)},
    title style={font=\scriptsize, yshift=-2pt},
    enlarge x limits=0.02,
    clip mode=individual,
]
    % Tren de pulsos: ancho tau=1 centrado en multiplos de T_0=4
    \addplot[utnblue, thick] coordinates {
        (-7,0) (-4.5,0) (-4.5,1) (-3.5,1) (-3.5,0)
        (-0.5,0) (-0.5,1) (0.5,1) (0.5,0)
        (3.5,0) (3.5,1) (4.5,1) (4.5,0) (7,0)
    };
    \node[anchor=south west, font=\footnotesize] at (axis cs:0.1, 1.0) {$1$};
    \node[font=\scriptsize] at (axis cs:-6.3, 0.5) {$\cdots$};
    \node[font=\scriptsize] at (axis cs:6.3, 0.5)  {$\cdots$};
\end{axis}

% ===== Panel (b): c_k vs k (indice discreto) =====
\begin{axis}[
    name=panel_b,
    at={($(panel_a.east)+(1.2cm,0)$)}, anchor=west,
    width=4.5cm, height=4cm,
    axis lines=center,
    axis line style={->, thick},
    xmin=-6.5, xmax=6.5, ymin=-0.10, ymax=0.38,
    xtick={-4,0,4},
    ytick={0, 0.25},
    yticklabels={$0$, $\tfrac{1}{4}$},
    every tick label/.style={font=\footnotesize},
    xlabel={$k$},
    xlabel style={anchor=west, at={(axis description cs:1.0,0.25)}},
    ylabel={$c_k$},
    ylabel style={font=\scriptsize, rotate=0, anchor=south east,
                  at={(axis description cs:0,1)}, xshift=-2pt},
    title={(b)},
    title style={font=\scriptsize, yshift=-2pt},
    enlarge x limits=0.02,
    clip mode=individual,
]
    \addplot+[ycomb, mark=*, mark size=0.5pt, utnblue, thick,
              mark options={fill=utnblue, draw=utnblue}]
        coordinates {
            (-6,0.0530) (-5,-0.0450) (-4,0) (-3,0.0750) (-2,0.1592)
            (-1,0.2251) (0,0.2500) (1,0.2251) (2,0.1592) (3,0.0750)
            (4,0) (5,-0.0450) (6,0.0530)
        };
\end{axis}

% ===== Panel (c): X(jw) impulsos en w = k*omega_0 =====
\begin{axis}[
    name=panel_c,
    at={($(panel_b.east)+(1.2cm,0)$)}, anchor=west,
    width=4.5cm, height=4cm,
    axis lines=center,
    axis line style={->, thick},
    xmin=-10.5, xmax=10.5, ymin=-0.55, ymax=2.25,
    xtick={-6.2832, 0, 6.2832},
    xticklabels={$-4\omega_0$,$0$,$4\omega_0$},
    ytick={0, 1.5708},
    yticklabels={$0$, $\tfrac{\pi}{2}$},
    every tick label/.style={font=\footnotesize},
    xlabel={$\omega$},
    xlabel style={anchor=west, at={(axis description cs:1.0,0.25)}},
    ylabel={$X(j\omega)$},
    ylabel style={font=\scriptsize, rotate=0, anchor=south east,
                  at={(axis description cs:0,1)}, xshift=-2pt},
    title={(c)},
    title style={font=\scriptsize, yshift=-2pt},
    enlarge x limits=0.02,
    clip mode=individual,
]
    % Impulsos como flechas verticales en omega = k*omega_0 = k*pi/2.
    % Altura = 2*pi*c_k.
    \draw[utnblue, thick, ->] (axis cs:0,0)         -- (axis cs:0,1.5708);
    \draw[utnblue, thick, ->] (axis cs:1.5708,0)    -- (axis cs:1.5708,1.4142);
    \draw[utnblue, thick, ->] (axis cs:-1.5708,0)   -- (axis cs:-1.5708,1.4142);
    \draw[utnblue, thick, ->] (axis cs:3.1416,0)    -- (axis cs:3.1416,1.0);
    \draw[utnblue, thick, ->] (axis cs:-3.1416,0)   -- (axis cs:-3.1416,1.0);
    \draw[utnblue, thick, ->] (axis cs:4.7124,0)    -- (axis cs:4.7124,0.4714);
    \draw[utnblue, thick, ->] (axis cs:-4.7124,0)   -- (axis cs:-4.7124,0.4714);
    % k=4 es cero, no se dibuja.
    \draw[utnblue, thick, ->] (axis cs:7.8540,0)    -- (axis cs:7.8540,-0.2828);
    \draw[utnblue, thick, ->] (axis cs:-7.8540,0)   -- (axis cs:-7.8540,-0.2828);
    \draw[utnblue, thick, ->] (axis cs:9.4248,0)    -- (axis cs:9.4248,0.3331);
    \draw[utnblue, thick, ->] (axis cs:-9.4248,0)   -- (axis cs:-9.4248,0.3331);
\end{axis}

\end{tikzpicture}
\end{document}
